\section{Часть 2. Методика оценки безопасности кода}
\label{sec:part2}

\textbf{Описание своего и предоставленного алгоритмов.}
Собственная реализация хранит секреты в ОЗУ только как шифртекст, затирает временный буфер открытого текста и отделяет секреты от обычных заметок на уровне типов и интерфейса.

Предоставленный алгоритм реализует тот же набор операций создания, обновления и удаления, однако поле секрета~\mdash  обычная строка без шифрования и без очистки буферов. Это типичная ошибка, при которой конфиденциальные данные в куче процесса неотличимы от любых других текстовых полей.

\textbf{Анализ безопасности предоставленного алгоритма.}
В предоставленном варианте секрет живёт в поле объекта столько, сколько существует запись. Валидация минимальна: проверяется лишь наличие полей, ограничения длины и затирание отсутствуют. Журнал не содержит тел секретов, однако этого недостаточно: при снятии дампа после создания записи открытый текст секрета стабильно читается как содержимое объекта. С точки зрения кода нет разделения <<хранить открыто / хранить защищённо>> и нет криптографической обработки сразу после ввода. В аварийной ситуации и при анализе памяти атакующий с доступом к дампу процесса получает секреты в явном виде. По процессорному времени вариант сопоставим с собственной реализацией, поскольку объём данных невелик; выигрыш по CPU не компенсирует потерю конфиденциальности.

\textbf{Описание методики.}
Методика оценивает не внешний вид интерфейса, а то, как код обращается с данными в ОЗУ. Этапы: инвентаризация сущностей и мест хранения; проверка реакции на пустой, чрезмерно длинный и некорректный ввод; анализ времени жизни открытого текста; проверка обновления и удаления по дампам; учёт аварийных сценариев (журнал, дамп процесса); оценка ресурсов; сравнение по чек-листу. Каждый критерий оценивается баллами 0 (нет), 1 (частично), 2 (да). Сумма 0--5 соответствует низкому уровню, 6--11~\mdash  среднему, 12--16~\mdash  высокому относительно учебных требований работы с памятью.

\begin{table}[ht]
	\begin{tabularx}{\textwidth}{|c|X|}
		\hline
		\textbf{ID} & \textbf{Критерий} \\
		\hline
		C1 & Конфиденциальные данные не хранятся открытым текстом дольше необходимого \\
		\hline
		C2 & После ввода применяется шифрование или необратимое хеширование \\
		\hline
		C3 & Есть затирание временных буферов открытого текста \\
		\hline
		C4 & Конфиденциальные и неконфиденциальные сущности разделены \\
		\hline
		C5 & Валидация обязательных полей и ограничений длины на сервере \\
		\hline
		C6 & Секреты не пишутся в журнал приложения \\
		\hline
		C7 & После удаления маркеры открытого текста секрета не находятся в дампе \\
		\hline
		C8 & Использование памяти и процессорного времени допускает наблюдение и оценку \\
		\hline
	\end{tabularx}
	\caption{Чек-лист методики оценки безопасности кода}
	\label{tab:checklist}
\end{table}

\pagebreak
\textbf{Сравнение по методике.}

\begin{table}[ht]
	\begin{tabularx}{\textwidth}{|l|c|c|X|}
		\hline
		\textbf{Критерий} & \textbf{Свой} & \textbf{Предоставленный} & \textbf{Комментарий} \\
		\hline
		C1 & 1 & 0 & у своего возможны краткие копии запроса \\
		\hline
		C2 & 2 & 0 & Fernet против открытого текста \\
		\hline
		C3 & 2 & 0 & затирание буфера только в своём варианте \\
		\hline
		C4 & 2 & 1 & у предоставленного тип секрета есть, семантика открытая \\
		\hline
		C5 & 2 & 1 & у своего есть ограничения длины \\
		\hline
		C6 & 2 & 2 & оба не журналируют секреты \\
		\hline
		C7 & 1 & 0 & после удаления у своего маркеры секретов не найдены \\
		\hline
		C8 & 2 & 1 & оба допускают снятие дампа и замер ресурсов \\
		\hline
		\textbf{Итого} & \textbf{14/16} & \textbf{5/16} & высокий / низкий уровень \\
		\hline
	\end{tabularx}
	\caption{Сравнение реализаций по чек-листу методики}
	\label{tab:cmp}
\end{table}

Методика различает варианты по существу лабораторной: предоставленный алгоритм проигрывает по шифрованию, затиранию буферов и следам в дампе; собственная реализация снижает риск чтения секрета из ОЗУ, хотя полностью устранить кратковременные копии строк в интерпретаторе нельзя.
